% ==========================================================================
% The Event Density Architecture:
% Foundations, Invariants, and Computational Verification
%
% Allen Proxmire
% Event Density Research Program, 2026
% ==========================================================================

\documentclass[11pt,twoside,openright]{book}

% --- Core mathematics ---
\usepackage{amsmath,amssymb,amsthm,mathtools}

% --- Layout ---
\usepackage[
  paperwidth=6.5in,
  paperheight=9.5in,
  inner=1in,
  outer=0.75in,
  top=0.85in,
  bottom=1in,
  headheight=14pt
]{geometry}

% --- Typography ---
\usepackage{microtype}
\usepackage{setspace}
\onehalfspacing

% --- Headers and footers ---
\usepackage{fancyhdr}
\pagestyle{fancy}
\fancyhf{}
\fancyhead[LE]{\small\itshape\nouppercase{\leftmark}}
\fancyhead[RO]{\small\itshape\nouppercase{\rightmark}}
\fancyfoot[C]{\thepage}
\renewcommand{\headrulewidth}{0.4pt}

% --- Section titles ---
\usepackage{titlesec}
\titleformat{\part}[display]
  {\centering\Huge\bfseries}
  {\partname~\thepart}{20pt}{\Huge}
\titleformat{\chapter}[display]
  {\Large\bfseries}
  {\chaptertitlename~\thechapter}{12pt}{\LARGE}
\titlespacing*{\chapter}{0pt}{-20pt}{30pt}

% --- References and links ---
\usepackage[
  colorlinks=true,
  linkcolor=blue!60!black,
  citecolor=green!50!black,
  urlcolor=blue!70!black
]{hyperref}
\usepackage[capitalise,noabbrev]{cleveref}

% --- Graphics ---
\usepackage{graphicx}
\usepackage{xcolor}
\usepackage{tikz}

% --- Lists ---
\usepackage{enumitem}
\setlist{nosep,leftmargin=*}

% --- Tables ---
\usepackage{booktabs}
\usepackage{array}
\usepackage{longtable}

% --- Bibliography ---
\usepackage[
  backend=biber,
  style=numeric,
  sorting=none,
  maxnames=4
]{biblatex}
\addbibresource{ed_references.bib}

% ==========================================================================
% Theorem environments
% ==========================================================================
\theoremstyle{definition}
\newtheorem{definition}{Definition}[chapter]

\theoremstyle{plain}
\newtheorem{proposition}{Proposition}[chapter]
\newtheorem{theorem}{Theorem}[chapter]
\newtheorem{lemma}{Lemma}[chapter]

\theoremstyle{definition}
\newtheorem{principle}{Principle}[chapter]
\newtheorem{prediction}{Prediction}[chapter]

\theoremstyle{remark}
\newtheorem{remark}{Remark}[chapter]

% ==========================================================================
% Custom commands
% ==========================================================================
\newcommand{\rhomax}{\rho_{\max}}
\newcommand{\rhostar}{\rho^{*}}
\newcommand{\Fop}{F[\rho]}
\newcommand{\Fbar}{\bar{F}[\rho]}
\newcommand{\Dgrad}{\mathcal{D}_{\mathrm{grad}}}
\newcommand{\Dpen}{\mathcal{D}_{\mathrm{pen}}}
\newcommand{\Dpart}{\mathcal{D}_{\mathrm{part}}}
\newcommand{\Dtotal}{\mathcal{D}_{\mathrm{total}}}
\newcommand{\Energy}{\mathcal{E}}
\newcommand{\Disc}{\mathscr{D}_0}
\newcommand{\UED}{\mathcal{U}_{\mathrm{ED}}}
\newcommand{\att}{\mathrm{att}}
\newcommand{\CV}{\mathrm{CV}}
\DeclareMathOperator{\corr}{corr}

% ==========================================================================
\begin{document}

% ==========================================================================
% FRONT MATTER
% ==========================================================================
\frontmatter

% --- Title page ---
\begin{titlepage}
\centering
\vspace*{2cm}
{\Huge\bfseries The Event Density Architecture}\\[0.8cm]
{\Large Foundations, Invariants, and Computational Verification}\\[2.5cm]
{\large\itshape Allen Proxmire}\\[0.5cm]
{\normalsize Event Density Research Program}\\[4cm]
{\normalsize 2026}
\vfill
\end{titlepage}

% --- Copyright ---
\thispagestyle{empty}
\vspace*{\fill}
\noindent\textcopyright\ 2026 Allen Proxmire.\\
Event Density Research Program.\\[1em]
\noindent ED-SIM v1.0.0.
\vspace*{\fill}
\clearpage

% --- Abstract ---
\chapter*{Abstract}
\addcontentsline{toc}{chapter}{Abstract}

We present ED-SIM~v1, a reproducible computational pipeline for the Event Density (ED) architectural ontology. The pipeline solves the canonical ED partial differential equation across a three-dimensional parameter space \((D, A, N_m)\) and computes sixteen families of attractor invariants that characterise the structural properties of the ED dynamical system. Three meta-analyses---parameter universality, cross-invariant consistency, and embedding collapse---synthesise the invariant families into a global architectural verdict, encoded in a machine-verifiable ED Consistency Certificate.

The principal findings are: (i)~the ED system possesses a unique point attractor \((\rhostar, 0)\) across all tested parameter regimes, confirmed by zero positive Lyapunov exponents and low effective attractor dimension; (ii)~the sixteen invariant families are mutually consistent and structurally invariant under parameter variation, with coefficient of variation below 5\% for the core families; (iii)~the three-stage convergence structure predicted by the analytic theory (global bounds, algebraic decay, exponential decay) is reproduced quantitatively in every admissible run; and (iv)~the complete pipeline---from raw PDE integration through invariant computation to certificate generation---is fully reproducible from a single command.

The invariant atlas constitutes the first systematic numerical verification of the ED architectural ontology and provides the empirical foundation for the physical predictions of the ED Applications Paper.

% --- Table of contents ---
\tableofcontents
\listoffigures
\listoftables

% --- Notation table ---
\chapter*{Notation}
\addcontentsline{toc}{chapter}{Notation}

\begin{longtable}{@{}>{$}l<{$}p{9cm}@{}}
\toprule
\textbf{Symbol} & \textbf{Meaning} \\
\midrule
\endhead
\rho(x,t) & Density field \\
v(t) & Participation variable \\
\Fop & Density operator \\
M(\rho) & Mobility function \\
P(\rho) & Penalty function \\
D & Direct-channel weight \\
H = 1 - D & Participation-channel weight \\
\zeta & Participation damping rate \\
\tau & Participation time scale \\
\rhostar & Penalty equilibrium density \\
\rhomax & Capacity bound \\
M_0, \beta & Mobility parameters \\
P_0 & Penalty slope \\
\Omega = [0, L] & Spatial domain \\
\varphi_k(x) & Neumann eigenfunctions \\
\mu_k = (k\pi/L)^2 & Eigenvalues \\
a_k(t) & Modal amplitudes \\
E_k = |a_k|^2 & Modal energies \\
\alpha_k = M_*\mu_k + P_*' & Modal decay coefficients \\
\Disc & Homogeneous-mode discriminant \\
\gamma_0 & Homogeneous-mode decay rate \\
\Energy[\rho, v] & Energy functional \\
\Phi(\rho) & Density potential \\
\Dgrad, \Dpen, \Dpart & Dissipation channels \\
\langle \cdot \rangle_\att & Mean over attractor window \\
\CV & Coefficient of variation \\
U & Universality score \\
C & Consistency score \\
D_{\mathrm{KY}} & Kaplan--Yorke dimension \\
D_{\mathrm{eff}} & Effective PCA dimension \\
\bottomrule
\end{longtable}


% --- Preface ---
\chapter*{Preface}
\addcontentsline{toc}{chapter}{Preface}

This monograph presents the Event Density (ED) architecture as a unified mathematical framework, develops its invariant theory, and documents its computational verification through the ED-SIM-01 pipeline. The work draws on three preceding documents---the Architectural Canon, the Rigour Paper, and the Applications Paper---and synthesises them with the Numerical Atlas, the Simulation Suite, and the Reproducibility Suite into a single, self-contained account.

The monograph is organised in six parts. \cref{part:foundations} establishes the foundations: the governing equations, the seven canonical principles, and the architectural predictions. \cref{part:invariants} develops the invariant theory: sixteen families of structural quantities that characterise the ED attractor from independent angles. \cref{part:computation} describes the computation: the ED-SIM-01 pipeline, its numerical methods, and its parameter-space coverage. \cref{part:synthesis} presents the synthesis: the invariant atlas, the global verdict, and the ED Architecture Certificate. \cref{part:reproducibility} addresses reproducibility. \cref{part:outlook} discusses applications, open problems, and future directions.

The reader is assumed to be familiar with partial differential equations, dynamical systems, and scientific computing at the level of a graduate course. No prior knowledge of the ED framework is required.


% ==========================================================================
% MAIN MATTER
% ==========================================================================
\mainmatter

% **************************************************************************
\part{Foundations}\label{part:foundations}
% **************************************************************************

% ------------------------------------------------------------------
\chapter{Introduction}\label{ch:intro}
% ------------------------------------------------------------------

\section{Motivation}\label{sec:motivation}

Physical theories typically operate at one of two levels. At the \emph{microscopic} level, they describe the fundamental constituents---particles, fields, quanta---and derive macroscopic behaviour from first principles. At the \emph{phenomenological} level, they model observed phenomena with effective equations whose parameters are fit to data. Between these extremes lies a third possibility: \emph{architectural} theories, which describe the structural organisation of dynamical systems without specifying their microscopic content.

Event Density is an architectural theory. It does not postulate new particles, new fields, or new forces. Instead, it identifies a set of structural principles---seven in number, each irreducible---that constrain how a density field can evolve, dissipate, and converge. The claim is that these principles are not specific to any one physical domain: they apply equally to quantum systems, galactic dynamics, condensed matter, and nonlinear optics, because they describe the \emph{architecture} of the dynamics, not the \emph{substrate}.

This claim is testable. The seven principles generate a specific partial differential equation---the canonical ED PDE---whose solutions can be computed, analysed, and compared to physical observations. The present monograph develops the mathematical infrastructure for this programme.

\section{The Central Question}\label{sec:central-question}

The central question of this monograph is:

\begin{quote}
\itshape Does the canonical ED PDE, solved without approximation across a systematic parameter space, exhibit the structural properties predicted by the seven principles?
\end{quote}

This is not a question about parameter fitting. The ED PDE has no free parameters beyond its canonical set \((D, \zeta, \tau, \rhostar, \rhomax)\), which specify the architecture, not the data. The question is whether the architecture itself---the seven principles, the operator structure, the constitutive constraints---produces the predicted qualitative behaviour: a unique point attractor, a three-stage convergence, a universal modal hierarchy, and a structurally invariant spectral fingerprint.

\section{Strategy}\label{sec:strategy}

The strategy is invariant analysis. We define sixteen families of scalar and vector quantities that capture different structural aspects of the long-time solution. We compute these invariants across sixty-four parameter regimes. We test whether they are invariant (constant across regimes), consistent (mutually compatible), and stable (robust to perturbation). The results are synthesised into a single verdict---the ED Architecture Certificate---which encodes the global structural consistency of the framework.

\section{Overview of the Monograph}\label{sec:overview}

\begin{table}[h]
\centering
\caption{Structure of the monograph.}\label{tab:overview}
\begin{tabular}{@{}lll@{}}
\toprule
Part & Chapters & Content \\
\midrule
I   & 1--5   & Foundations: equations, principles, predictions \\
II  & 6--9   & Invariants: definitions, interpretations, dependencies \\
III & 10--14 & Computation: pipeline, methods, extraction, meta-analyses \\
IV  & 15--18 & Synthesis: atlas, verdict, certificate, interpretation \\
V   & 19--21 & Reproducibility: suite, validation, extension \\
VI  & 22--24 & Outlook: applications, open problems, conclusion \\
\bottomrule
\end{tabular}
\end{table}


% ------------------------------------------------------------------
\chapter{The Event Density Framework}\label{ch:framework}
% ------------------------------------------------------------------

\section{Conceptual Picture}\label{sec:conceptual}

The ED framework describes the evolution of a \emph{density field} \(\rho(x,t)\) on a bounded spatial domain~\(\Omega\). The density represents a generic conserved or semi-conserved quantity---particle density, energy density, information density, or any scalar field whose dynamics are governed by transport, reaction, and feedback.

The evolution is driven by three competing mechanisms:

\begin{enumerate}
\item \textbf{Diffusion.} The density spreads from regions of high concentration to regions of low concentration. The diffusion rate is governed by a \emph{mobility function} \(M(\rho)\) that depends on the local density.
\item \textbf{Penalty.} The density is pulled toward a preferred equilibrium value \(\rhostar\) by a restoring force. The penalty function \(P(\rho)\) vanishes at \(\rhostar\) and grows monotonically away from it.
\item \textbf{Feedback.} The system integrates a spatial average of its own operator output and feeds it back into the density equation through a \emph{participation variable}~\(v(t)\). This creates a global coupling: the evolution at each point depends on the state of the entire domain.
\end{enumerate}

The interplay of these three mechanisms produces a rich dynamical landscape---oscillatory, monotonic, or critical, depending on the parameters---but the long-time outcome is always the same: the density converges to~\(\rhostar\), the gradients vanish, and the participation decays to zero.

\section{The Density Field}\label{sec:density-field}

The density \(\rho(x,t)\) is a scalar function of position \(x \in \Omega\) and time \(t \geq 0\). It takes values in the open interval \((0, \rhomax)\), where \(\rhomax > 0\) is the \emph{capacity bound}---the maximum density that the system can sustain. The boundary \(\rho = 0\) represents the vacuum; the boundary \(\rho = \rhomax\) represents the capacity horizon.

\begin{definition}[Admissible state]\label{def:admissible}
A density profile \(\rho(\cdot, t)\) is \emph{admissible} if \(\rho(x,t) \in (0, \rhomax)\) for all \(x \in \Omega\) and \(\rho(\cdot, t) \in H^1(\Omega)\).
\end{definition}

The admissible region \(\mathcal{O} = \{\rho \in H^1(\Omega) : 0 < \rho < \rhomax\}\) is an open subset of the Sobolev space \(H^1(\Omega)\). The ED dynamics preserves admissibility: if the initial condition is admissible, the solution remains admissible for all time~\cite{rigour}.

\section{The Participation Variable}\label{sec:participation}

The participation variable \(v(t)\) is a scalar that integrates the spatial average of the density operator and feeds it back into the evolution equation. It is governed by an ordinary differential equation with exponential damping:
\begin{equation}\label{eq:participation}
\dot{v} = \frac{1}{\tau}\bigl(\Fbar - \zeta\,v\bigr),
\end{equation}
where \(\Fbar = |\Omega|^{-1}\int_\Omega \Fop\,dx\) is the spatial average of the density operator, \(\tau > 0\) is the integration time scale, and \(\zeta > 0\) is the damping rate.

The participation variable creates a feedback loop: the density generates \(\Fop\), which is averaged and integrated into~\(v\), which in turn modifies the density evolution. This loop is the mechanism by which global information (the spatial average) influences local dynamics (the pointwise evolution).

\section{The Equilibrium}\label{sec:equilibrium}

\begin{definition}[Canonical equilibrium]\label{def:equilibrium}
The canonical equilibrium of the ED system is the pair \((\rhostar, 0)\), where \(\rho(x) = \rhostar\) is the spatially uniform density at the penalty zero and \(v = 0\) is the quiescent participation.
\end{definition}

At equilibrium, \(\nabla\rho = 0\), so \(M(\rhostar)\nabla^2\rho = 0\) and \(M'(\rhostar)|\nabla\rho|^2 = 0\). Since \(P(\rhostar) = 0\), the operator gives \(F[\rhostar] = 0\), which implies \(\dot{v} = -\zeta v/\tau\) and therefore \(v \to 0\) exponentially. The equilibrium is self-consistent.


% ------------------------------------------------------------------
\chapter{Governing Equations and Operators}\label{ch:equations}
% ------------------------------------------------------------------

\section{The Canonical ED PDE}\label{sec:canonical-pde}

The canonical ED system is the coupled PDE--ODE:
\begin{subequations}\label{eq:canonical}
\begin{align}
\partial_t\rho &= D\,\Fop + H\,v, \label{eq:canonical-a}\\
\dot{v} &= \frac{1}{\tau}\bigl(\Fbar - \zeta\,v\bigr), \label{eq:canonical-b}
\end{align}
\end{subequations}
on the domain \(\Omega = [0,L]\) with Neumann boundary conditions \(\partial_x\rho(0,t) = \partial_x\rho(L,t) = 0\) and initial conditions \(\rho(x,0) = \rho_0(x)\), \(v(0) = v_0\).

\section{The Density Operator}\label{sec:operator}

The density operator \(\Fop\) is defined by
\begin{equation}\label{eq:operator}
\Fop = M(\rho)\,\nabla^2\rho + M'(\rho)\,|\nabla\rho|^2 - P(\rho).
\end{equation}
This operator has three terms:
\begin{enumerate}
\item \textbf{Linear diffusion}: \(M(\rho)\,\nabla^2\rho\). A second-order parabolic term with density-dependent coefficient~\(M(\rho)\).
\item \textbf{Quadratic gradient}: \(M'(\rho)\,|\nabla\rho|^2\). A nonlinear term that couples spatial modes through the triad selection rule.
\item \textbf{Penalty reaction}: \(-P(\rho)\). A zeroth-order restoring term that drives \(\rho\) toward~\(\rhostar\).
\end{enumerate}

\begin{remark}\label{rem:divergence}
The operator \(\Fop\) can be written in divergence form as \(\Fop = \nabla \cdot (M(\rho)\nabla\rho) - P(\rho)\), which makes the mass-flux structure explicit: the density flux is \(J = -M(\rho)\nabla\rho\), and the penalty acts as a distributed source/sink.
\end{remark}

\section{Constitutive Functions}\label{sec:constitutive}

\begin{definition}[Canonical constitutive pair]\label{def:constitutive}
The canonical mobility and penalty functions are:
\begin{equation}\label{eq:constitutive}
M(\rho) = M_0\,(\rhomax - \rho)^\beta, \qquad P(\rho) = P_0\,(\rho - \rhostar),
\end{equation}
where \(M_0 > 0\), \(\beta > 0\), \(P_0 > 0\), and \(0 < \rhostar < \rhomax\).
\end{definition}

The mobility \(M(\rho)\) is positive on \((0, \rhomax)\) and vanishes at \(\rho = \rhomax\). This vanishing is the \emph{mobility collapse}---the mechanism that enforces the capacity bound. As \(\rho \to \rhomax\), the diffusion rate goes to zero, creating an impenetrable barrier.

The penalty \(P(\rho)\) vanishes at \(\rho = \rhostar\) and is positive for \(\rho > \rhostar\), negative for \(\rho < \rhostar\). It is the restoring force that drives the density toward equilibrium.

\section{Channel Structure}\label{sec:channels}

The density equation~\eqref{eq:canonical-a} is a weighted sum of two channels:
\begin{itemize}
\item The \textbf{direct channel}: \(D\,\Fop\), with weight \(D \in (0,1)\).
\item The \textbf{participation channel}: \(H\,v\), with weight \(H = 1 - D\).
\end{itemize}
The channel complementarity condition \(D + H = 1\) ensures that the two channels partition the operator output without amplification. The parameter~\(D\) controls the balance: at \(D = 1\), the system is purely direct (no feedback); at \(D = 0\), it is purely mediated (all feedback). The canonical range is \(D \in (0,1)\).

\section{The Energy Functional}\label{sec:energy}

\begin{definition}[Energy functional]\label{def:energy}
The canonical energy functional is:
\begin{equation}\label{eq:energy}
\Energy[\rho, v] = \int_\Omega \Phi(\rho)\,dx + \frac{\tau H}{2}\,v^2,
\end{equation}
where the density potential~\(\Phi\) is defined by
\begin{equation}\label{eq:potential}
\Phi(\rho) = \int_{\rhostar}^{\rho} \frac{P(s)}{M(s)}\,ds.
\end{equation}
\end{definition}

\begin{proposition}[Energy dissipation]\label{prop:dissipation}
Along solutions of~\eqref{eq:canonical}, the energy functional satisfies
\begin{equation}\label{eq:dissipation}
\frac{d\Energy}{dt} = -\Dgrad - \Dpen - \Dpart \leq 0,
\end{equation}
where the three dissipation channels are:
\begin{align*}
\Dgrad &= D\int_\Omega \frac{|\nabla\rho|^2}{M(\rho)}\cdot P'(\rho)\,dx, \\
\Dpen &= D\int_\Omega \frac{P(\rho)^2}{M(\rho)}\,dx, \\
\Dpart &= \frac{H\zeta}{\tau}\,v^2.
\end{align*}
\end{proposition}

The energy is a Lyapunov functional: it decreases monotonically and is bounded below by zero. This is the mathematical foundation of the three-stage convergence.

\section{Spectral Structure}\label{sec:spectral}

Expanding the perturbation \(u(x,t) = \rho(x,t) - \rhostar\) in the Neumann eigenbasis:
\begin{equation}\label{eq:eigenbasis}
u(x,t) = \sum_{k=0}^{\infty} a_k(t)\,\varphi_k(x), \qquad \varphi_k(x) = \sqrt{\frac{2}{L}}\cos\!\Bigl(\frac{k\pi x}{L}\Bigr),
\end{equation}
with eigenvalues \(\mu_k = (k\pi/L)^2\), the linearised dynamics of the \(k\)-th mode has decay rate
\begin{equation}\label{eq:decay-rate}
\alpha_k = M_*\,\mu_k + P_*',
\end{equation}
where \(M_* = M(\rhostar)\) and \(P_*' = P'(\rhostar)\). Since \(M_* > 0\) and \(P_*' > 0\), the decay rates are positive and increase with~\(k\): higher modes decay faster. This is the \emph{modal hierarchy}.


% ------------------------------------------------------------------
\chapter{The Seven Principles of ED}\label{ch:principles}
% ------------------------------------------------------------------

The ED architecture is built on seven irreducible principles. Each principle constrains the structure of the governing equation and excludes alternatives. Together, they determine the qualitative dynamics uniquely.

\begin{principle}[Operator Structure]\label{prin:operator}
The density operator \(\Fop\) consists of exactly three terms: mobility-weighted diffusion \(M(\rho)\nabla^2\rho\), gradient-squared nonlinearity \(M'(\rho)|\nabla\rho|^2\), and penalty reaction~\(-P(\rho)\). No other terms are present.
\end{principle}

\begin{remark}\label{rem:prin1}
This excludes advection, higher-order diffusion, nonlocal operators, and stochastic forcing from the canonical system. The three terms are the minimal set that produces a well-posed parabolic PDE with a unique equilibrium and nonlinear mode coupling.
\end{remark}

\begin{principle}[Channel Complementarity]\label{prin:channels}
The density evolution is a weighted sum of a direct channel \(D\,\Fop\) and a participation channel \(H\,v\), with \(D + H = 1\) and \(D \in (0,1)\).
\end{principle}

\begin{remark}\label{rem:prin2}
The complementarity condition prevents unbounded amplification and ensures that the total operator weight is unity. It creates a trade-off: increasing the direct channel decreases the feedback, and vice versa.
\end{remark}

\begin{principle}[Penalty Equilibrium]\label{prin:penalty}
The penalty function satisfies \(P(\rhostar) = 0\) and \(P'(\rhostar) > 0\). The equilibrium~\(\rhostar\) is unique.
\end{principle}

\begin{remark}\label{rem:prin3}
The strict positivity of \(P'(\rhostar)\) ensures that the penalty is a genuine restoring force, not a neutral fixed point. Combined with \cref{prin:operator,prin:capacity}, it guarantees that \((\rhostar, 0)\) is a global attractor.
\end{remark}

\begin{principle}[Mobility Capacity Bound]\label{prin:capacity}
The mobility satisfies \(M(\rhomax) = 0\) and \(M(\rho) > 0\) for \(\rho \in (0, \rhomax)\).
\end{principle}

\begin{remark}\label{rem:prin4}
The vanishing of~\(M\) at \(\rhomax\) creates an impassable barrier: the density cannot reach or exceed~\(\rhomax\). This is the \emph{mobility horizon}. Near \(\rhomax\), the mobility collapse freezes the dynamics, creating an effective boundary layer.
\end{remark}

\begin{principle}[Participation Feedback]\label{prin:feedback}
The participation variable \(v(t)\) integrates the spatial average of \(\Fop\) with time scale~\(\tau\) and damping rate~\(\zeta\), and feeds back into the density equation through the channel weight~\(H\).
\end{principle}

\begin{remark}\label{rem:prin5}
The participation loop is the mechanism by which global information enters the local dynamics. It introduces a second time scale~(\(\tau\)) and a second damping parameter~(\(\zeta\)) into the system.
\end{remark}

\begin{principle}[Damping Discriminant]\label{prin:discriminant}
The linearised dynamics of the homogeneous mode is governed by the discriminant
\begin{equation}\label{eq:discriminant}
\Disc = \bigl(D\,P_*' - \zeta/\tau\bigr)^2 - 4\,H\,P_*'/\tau,
\end{equation}
which classifies the dynamics as oscillatory (\(\Disc < 0\)), critical (\(\Disc = 0\)), or monotonic (\(\Disc > 0\)).
\end{principle}

\begin{remark}\label{rem:prin6}
The discriminant is the organising parameter of the regime geometry. It determines whether the approach to equilibrium involves oscillation (spiral sheet), pure decay (monotonic cone), or transitional dynamics (critical surface).
\end{remark}

\begin{principle}[Nonlinear Triad Coupling]\label{prin:triad}
The quadratic term \(M'(\rho)|\nabla\rho|^2\) generates inter-modal coupling with the selection rule: modes~\(m\) and~\(n\) interact to produce modes \(|m-n|\) and \(m+n\).
\end{principle}

\begin{remark}\label{rem:prin7}
The triad selection rule is a consequence of the product structure of the gradient-squared term in the Neumann eigenbasis. It determines which spectral modes can exchange energy and controls the shape of the nonlinear cascade.
\end{remark}


% ------------------------------------------------------------------
\chapter{Architectural Predictions}\label{ch:predictions}
% ------------------------------------------------------------------

The seven principles, combined with the analytic theory of the Rigour Paper~\cite{rigour}, generate the following structural predictions:

\begin{prediction}[Unique Point Attractor]\label{pred:attractor}
The canonical equilibrium \((\rhostar, 0)\) is a global attractor: every admissible initial condition converges to \((\rhostar, 0)\) as \(t \to \infty\).
\end{prediction}

\begin{prediction}[Three-Stage Convergence]\label{pred:three-stage}
The convergence proceeds through three stages:
\begin{enumerate}
\item Stage~I (Global Bounds): the energy functional is uniformly bounded and monotonically decreasing.
\item Stage~II (Algebraic Decay): gradients and participation decay algebraically via Barbalat's lemma.
\item Stage~III (Exponential Decay): the solution enters the local stability basin and converges exponentially at a rate determined by the spectral gap.
\end{enumerate}
\end{prediction}

\begin{prediction}[Non-Positive Lyapunov Spectrum]\label{pred:lyapunov}
All Lyapunov exponents of the attractor are non-positive. The Kaplan--Yorke dimension is zero.
\end{prediction}

\begin{prediction}[Modal Hierarchy]\label{pred:hierarchy}
The decay rate of the \(k\)-th eigenmode is \(\sigma_k = D\,\alpha_k\), where \(\alpha_k = M_*\mu_k + P_*'\) increases monotonically with~\(k\). Higher modes decay faster.
\end{prediction}

\begin{prediction}[Triad Selection Rule]\label{pred:triad}
Modes~\(m\) and~\(n\) interact only to produce modes \(|m-n|\) and \(m+n\). No other mode coupling exists.
\end{prediction}

\begin{prediction}[Universal Dissipation Partition]\label{pred:dissipation}
The three dissipation channels \((\Dgrad, \Dpen, \Dpart)\) converge to a fixed partition that depends only on the canonical parameters, not on the initial condition.
\end{prediction}

\begin{prediction}[Structural Invariance]\label{pred:invariance}
The qualitative attractor structure---modal hierarchy, triad rule, dissipation partition, spectral fingerprint---is invariant under parameter variation within the universality class~\(\UED\).
\end{prediction}

These seven predictions are the targets of the invariant atlas. Each is tested numerically by one or more of the sixteen invariant families defined in \cref{part:invariants}.


% **************************************************************************
\part{Invariants}\label{part:invariants}
% **************************************************************************

% ------------------------------------------------------------------
\chapter{Why Invariants}\label{ch:why-invariants}
% ------------------------------------------------------------------

\section{The Concept}\label{sec:invariant-concept}

An \emph{invariant} of the ED system is a quantity computed from the long-time solution that is approximately constant across parameter regimes. The word ``approximately'' is quantified by the coefficient of variation:
\begin{equation}\label{eq:cv}
\CV(Q) = \frac{\sigma(Q)}{\mu(Q)},
\end{equation}
where \(\mu(Q)\) and \(\sigma(Q)\) are the mean and standard deviation of the quantity~\(Q\) across all admissible runs. The invariance verdict is:

\begin{table}[h]
\centering
\caption{Invariance verdict thresholds.}\label{tab:cv-thresholds}
\begin{tabular}{@{}ll@{}}
\toprule
\(\CV\) & Verdict \\
\midrule
\(< 0.05\) & \textsc{invariant} \\
\(< 0.15\) & \textsc{weakly invariant} \\
\(\geq 0.15\) & \textsc{not invariant} \\
\bottomrule
\end{tabular}
\end{table}

\section{Why Sixteen Families}\label{sec:sixteen}

A single invariant proves nothing: it could be invariant by coincidence. Sixteen invariants probing the attractor from independent angles---spectral, dynamical, topological, and statistical---provide a structural test. Their mutual consistency (measured by the cross-invariant correlation matrix) is a stronger claim than any individual invariance.

\section{The Attractor Window}\label{sec:attractor-window}

All invariants are computed over the \emph{attractor window}: the last fraction of the time series, where the solution has entered (or nearly entered) the exponential-decay regime.

\begin{definition}[Attractor window]\label{def:attractor-window}
Let~\(T\) be the final integration time. The attractor window is the interval \([T_\att, T]\) where \(T_\att = (1-f)\,T\) and~\(f\) is the window fraction. The default values are \(f = 0.10\) for late-time averages and \(f = 0.20\) for convergence-rate fits.
\end{definition}


% ------------------------------------------------------------------
\chapter{The Sixteen Invariant Families}\label{ch:families}
% ------------------------------------------------------------------

This chapter defines each invariant family. For compactness, we use the notation \(\langle Q \rangle_\att\) for the mean of \(Q(t)\) over the attractor window.

\section{Family 1: Low-Mode Collapse}\label{sec:fam-lowmode}

\begin{definition}\label{def:lowmode}
For each mode \(k = 0, \ldots, 5\), the late-time amplitude is:
\begin{equation}\label{eq:lowmode}
m_k^* = \bigl\langle |a_k(t)| \bigr\rangle_\att.
\end{equation}
\end{definition}

\noindent\emph{Tests:} \cref{pred:attractor} (point attractor). If \((\rhostar, 0)\) is the unique attractor, then \(m_k^* \to 0\) for all~\(k\).

\section{Family 2: Mode-Energy Ratios}\label{sec:fam-moderatio}

\begin{definition}\label{def:moderatio}
For each mode \(k = 1, \ldots, K\), the energy ratio is:
\begin{equation}\label{eq:moderatio}
R_k^* = \bigl\langle E_k(t) / E_{\mathrm{total}}(t) \bigr\rangle_\att, \qquad E_k = |a_k|^2.
\end{equation}
\end{definition}

\noindent\emph{Tests:} \cref{pred:hierarchy} (modal hierarchy). The energy distribution should converge to a fixed profile with \(R_1^* > R_2^* > \cdots\).

\section{Family 3: Spectral Complexity}\label{sec:fam-renyi}

\begin{definition}\label{def:renyi}
Let \(p_k(t) = E_k(t)/\sum_j E_j(t)\). The R\'{e}nyi entropy of order~\(q\) is:
\begin{equation}\label{eq:renyi}
H_q(t) = \frac{1}{1-q}\log\!\Bigl(\sum_k p_k(t)^q\Bigr), \qquad q \in \{0, 0.5, 1, 2, 3, 4\}.
\end{equation}
The invariant is \(H_q^* = \langle H_q(t) \rangle_\att\).
\end{definition}

\noindent\emph{Tests:} \cref{pred:hierarchy}. The multi-scale entropy profile characterises the spectral shape of the attractor.

\section{Family 4: Dissipation Partitions}\label{sec:fam-dissipation}

\begin{definition}\label{def:dissipation-part}
The dissipation ratios are:
\begin{equation}\label{eq:dissip-ratios}
R_{\mathrm{grad}}^* = \bigl\langle \Dgrad/\Dtotal \bigr\rangle_\att, \quad
R_{\mathrm{pen}}^* = \bigl\langle \Dpen/\Dtotal \bigr\rangle_\att, \quad
R_{\mathrm{part}}^* = \bigl\langle \Dpart/\Dtotal \bigr\rangle_\att.
\end{equation}
\end{definition}

\noindent\emph{Tests:} \cref{pred:dissipation}. The three-channel structure converges to a fixed partition.

\section{Family 5: Energy--Entropy Geometry}\label{sec:fam-EH}

\begin{definition}\label{def:EH}
The attractor point in the energy--entropy plane is:
\begin{equation}\label{eq:EH}
(E^*, H^*) = \bigl(\langle E(t) \rangle_\att,\; \langle H_1(t) \rangle_\att\bigr).
\end{equation}
\end{definition}

\noindent\emph{Tests:} \cref{pred:attractor,pred:dissipation}. All runs should converge to the same \((E^*, H^*)\).

\begin{figure}[ht]
\centering
% Placeholder for figure
\fbox{\parbox{0.8\textwidth}{\centering\vspace{3cm}\textit{Energy--entropy trajectory: parametric curve \((E(t), H(t))\) spiralling toward \((E^*, H^*)\).}\vspace{3cm}}}
\caption{Energy--entropy trajectory for the representative run (largest~\(N_m\)), coloured by time. The trajectory spirals inward and converges to the attractor point~\((E^*, H^*)\), marked with a star.}
\label{fig:EH-trajectory}
\end{figure}

\section{Family 6: Broadband Cascade}\label{sec:fam-cascade}

\begin{definition}\label{def:cascade}
Define logarithmic mode bins (\(k = 1\text{--}2, 3\text{--}4, 5\text{--}8, 9\text{--}16, 17\text{--}32\)) and the bin-energy ratios:
\begin{equation}\label{eq:cascade}
R_b^* = \bigl\langle B_b(t) / E_{\mathrm{total}}(t) \bigr\rangle_\att, \qquad B_b = \sum_{k \in \mathrm{bin}\,b} E_k.
\end{equation}
\end{definition}

\noindent\emph{Tests:} \cref{pred:hierarchy,pred:triad}. The cascade profile is shaped by the modal hierarchy and the triad selection rule.

\section{Family 7: Convergence Stability}\label{sec:fam-convergence}

\begin{definition}\label{def:convergence}
For each error signal \(e(t) \in \{|E(t) - E^*|, |H(t) - H^*|, |\Dtotal(t) - \mathcal{D}^*|\}\), identify the three convergence stages and fit the Stage~III exponential rate:
\begin{equation}\label{eq:convergence}
e(t) \sim C\,e^{-\sigma\,t}.
\end{equation}
\end{definition}

\noindent\emph{Tests:} \cref{pred:three-stage} (three-stage convergence).

\section{Family 8: Modal Correlations}\label{sec:fam-corr}

\begin{definition}\label{def:modal-corr}
The correlation matrix of modal energies over the attractor window is:
\begin{equation}\label{eq:modal-corr}
C_{ij} = \corr\bigl(E_i(t),\, E_j(t)\bigr), \qquad i, j = 1, \ldots, K.
\end{equation}
Summary invariants: mean off-diagonal correlation, spectral radius, condition number.
\end{definition}

\noindent\emph{Tests:} \cref{pred:triad}. Triad-connected modes should be correlated; unconnected modes should be independent.

\section{Family 9: Modal Overlap}\label{sec:fam-overlap}

\begin{definition}\label{def:overlap}
The nearest-neighbour overlap ratio for mode~\(k\) is:
\begin{equation}\label{eq:overlap}
O_k^* = \bigl\langle (E_{k-1} + E_k + E_{k+1}) / E_{\mathrm{total}} \bigr\rangle_\att.
\end{equation}
\end{definition}

\noindent\emph{Tests:} \cref{pred:hierarchy}. The overlap profile characterises the local spectral structure.

\section{Family 10: Phase Dynamics}\label{sec:fam-phase}

\begin{definition}\label{def:phase}
Extract modal phases \(\phi_k(t) = \arg(a_k(t))\). Compute:
\begin{itemize}
\item Phase drift rates \(d\phi_k/dt\) (linear fit over attractor window).
\item Phase coherence \(C_{ij}^\phi = |\langle e^{i(\phi_i - \phi_j)} \rangle|\).
\item Triad phase closure \(\Delta_{ijk} = \phi_i + \phi_j - \phi_k\) for triads \((i,j,k)\) with \(i + j = k\).
\end{itemize}
\end{definition}

\noindent\emph{Tests:} \cref{pred:triad,pred:dissipation}. Phase-locking in the oscillatory regime is a signature of triad coupling.

\section{Family 11: Phase--Amplitude Coupling}\label{sec:fam-pac}

\begin{definition}\label{def:pac}
The self-PAC for mode~\(k\) is:
\begin{equation}\label{eq:pac}
\rho_k = \corr\bigl(|a_k(t)|,\, \phi_k(t)\bigr).
\end{equation}
The cross-mode PAC matrix is \(\mathrm{PAC}_{ij} = \corr(|a_i|, \phi_j)\).
\end{definition}

\noindent\emph{Tests:} \cref{pred:attractor}. At the point attractor, phase and amplitude should decouple (\(|\rho_k| \approx 0\)).

\section{Family 12: Lyapunov Spectrum}\label{sec:fam-lyap}

\begin{definition}\label{def:lyap}
Construct the state vector \(\mathbf{x}(t) = [\operatorname{Re}(a_1), \operatorname{Im}(a_1), \ldots]\), compute finite-difference tangent vectors, and extract the Lyapunov exponents~\(\lambda_i\) via QR accumulation. The Kaplan--Yorke dimension is:
\begin{equation}\label{eq:KY}
D_{\mathrm{KY}} = j + \frac{\sum_{i=1}^j \lambda_i}{|\lambda_{j+1}|},
\end{equation}
where~\(j\) is the largest index such that \(\sum_{i=1}^j \lambda_i \geq 0\).
\end{definition}

\noindent\emph{Tests:} \cref{pred:lyapunov}. All \(\lambda_i \leq 0\) and \(D_{\mathrm{KY}} \approx 0\).

\begin{figure}[ht]
\centering
% Placeholder for figure
\fbox{\parbox{0.8\textwidth}{\centering\vspace{3cm}\textit{Lyapunov spectrum: \(\lambda_i\) vs index~\(i\), all non-positive.}\vspace{3cm}}}
\caption{Lyapunov spectrum. Left: \(\lambda_i\) versus index~\(i\) for the representative run, showing all exponents non-positive. Right: \(D_{\mathrm{KY}}\) versus~\(D\) for all runs, showing near-zero values.}
\label{fig:lyapunov}
\end{figure}

\section{Family 13: Attractor Manifold}\label{sec:fam-manifold}

\begin{definition}\label{def:manifold}
Perform PCA on the state vectors in the attractor window. The effective dimension is:
\begin{equation}\label{eq:Deff}
D_{\mathrm{eff}} = \min\!\Bigl\{m : \sum_{i=1}^m \lambda_i \geq 0.99\sum_j \lambda_j\Bigr\}.
\end{equation}
The spectral gap is \(\lambda_1/\lambda_2\) and the curvature proxy is \(\kappa = \lambda_2/\lambda_1\).
\end{definition}

\noindent\emph{Tests:} \cref{pred:attractor,pred:invariance}. A point attractor has \(D_{\mathrm{eff}} \leq 2\).

\section{Family 14: Perturbed-Attractor Stability}\label{sec:fam-perturbed}

\begin{definition}\label{def:perturbed}
Perturb the attractor state by \(\epsilon\,\boldsymbol{\xi}\) for \(\epsilon \in \{10^{-6}, 10^{-5}, 10^{-4}, 10^{-3}\}\) and fit the recovery rate:
\begin{equation}\label{eq:recovery}
\sum_k |a_k(t) - a_k^*|^2 \sim C\,e^{-\sigma_\epsilon\,t}.
\end{equation}
\end{definition}

\noindent\emph{Tests:} \cref{pred:attractor}. The recovery rate~\(\sigma_\epsilon\) should be independent of~\(\epsilon\) (linear stability).

\section{Family 15: Parameter Universality}\label{sec:fam-universality}

\begin{definition}\label{def:universality}
Construct the unified invariant vector~\(\mathbf{I}\) for each run. Standardise across runs. The universality score is:
\begin{equation}\label{eq:U-score}
U = \frac{1}{1 + \CV(\text{pairwise distances})}.
\end{equation}
\end{definition}

\noindent\emph{Tests:} \cref{pred:invariance}. \(U \approx 1\) confirms structural invariance.

\section{Family 16: Cross-Invariant Consistency}\label{sec:fam-consistency}

\begin{definition}\label{def:consistency}
For each pair of invariant families \((A, B)\), compute the mean pairwise correlation. The consistency score is:
\begin{equation}\label{eq:C-score}
C = \bigl\langle |\corr(A, B)| \bigr\rangle_{A \neq B}.
\end{equation}
\end{definition}

\noindent\emph{Tests:} \cref{pred:attractor,pred:three-stage,pred:lyapunov,pred:hierarchy,pred:triad,pred:dissipation,pred:invariance} collectively. \(C > 0.8\) confirms that the families describe a single coherent structure.


% ------------------------------------------------------------------
\chapter{Structural Interpretation of Each Invariant}\label{ch:interpretation}
% ------------------------------------------------------------------

\section{Spectral Invariants (Families 1--3, 6, 9)}

The spectral invariants characterise the \emph{shape} of the attractor in mode space. The low-mode amplitudes (Family~1) test whether any modes survive; the energy ratios (Family~2) test the distribution; the R\'{e}nyi entropies (Family~3) test the disorder; the broadband cascade (Family~6) tests the scale distribution; and the modal overlap (Family~9) tests the local spectral coupling.

Together, these five families provide a complete spectral fingerprint of the attractor. If the fingerprint is invariant across parameter regimes, the spectral shape is a structural property of the architecture, not a parameter-dependent artifact.

\section{Dynamical Invariants (Families 7, 10--11, 14)}

The dynamical invariants characterise the \emph{temporal} structure of the approach to equilibrium. The convergence stability (Family~7) tests the three-stage pattern; the phase dynamics (Family~10) test the oscillatory structure; the phase--amplitude coupling (Family~11) tests the nonlinear interactions; and the perturbation recovery (Family~14) tests the basin stability.

\section{Topological Invariants (Families 12--13)}

The topological invariants characterise the \emph{geometry} of the attractor. The Lyapunov spectrum (Family~12) measures the attractor's stability and dimension; the PCA manifold (Family~13) measures its effective geometry. Together, they confirm that the attractor is a point (zero-dimensional, non-positive Lyapunov spectrum) rather than a limit cycle, torus, or strange attractor.

\section{Statistical Invariants (Families 4--5, 8, 15--16)}

The statistical invariants characterise the \emph{macroscopic} properties: how energy is lost (dissipation partitions, Family~4), how energy and entropy relate (Family~5), how modes are correlated (Family~8), whether the structure is universal (Family~15), and whether the families agree (Family~16).


% ------------------------------------------------------------------
\chapter{Cross-Family Dependencies and Consistency}\label{ch:cross-family}
% ------------------------------------------------------------------

\section{Dependency Graph}

The sixteen families are not fully independent. Several share input data (modal amplitudes), and several test overlapping aspects of the attractor. The dependency structure is:
\begin{itemize}
\item Families 1--3, 6, 8--14 all depend on the modal amplitudes~\(a_k(t)\).
\item Families 4--5, 7 depend on the energy and dissipation time series.
\item Family~15 depends on all other families (it aggregates their summaries).
\item Family~16 depends on all other families (it correlates them).
\end{itemize}

\section{Expected Correlations}

Some cross-family correlations are structurally mandated:
\begin{itemize}
\item Low-mode collapse (Family~1) and mode-energy ratios (Family~2) are directly related.
\item Lyapunov spectrum (Family~12) and attractor manifold (Family~13) both measure the attractor's geometry, but from different perspectives (dynamical vs statistical). They should agree: \(D_{\mathrm{KY}} \approx 0\) implies \(D_{\mathrm{eff}} \leq 2\).
\item Dissipation partitions (Family~4) and convergence stability (Family~7) both depend on the energy dynamics, but the partitions measure the \emph{channel balance} while the convergence measures the \emph{rate}.
\end{itemize}

\section{The Consistency Score}

The cross-invariant consistency score~\(C\) (\cref{def:consistency}) quantifies the overall agreement. A high~\(C\) does not mean the families are redundant---it means they describe the same underlying structure from different angles. A low~\(C\) would indicate that different aspects of the attractor behave differently under parameter variation, which would challenge the architectural claim.

\section{Invariant Summary}

\begin{table}[ht]
\centering
\caption{Summary of the sixteen invariant families.}\label{tab:invariants}
\small
\begin{tabular}{@{}clllc@{}}
\toprule
\# & Family & Type & Key invariant & Predictions \\
\midrule
1  & Low-mode collapse       & Spectral    & \(m_k^*\)                                                 & 5.1 \\
2  & Mode-energy ratios      & Spectral    & \(R_k^*\)                                                 & 5.4 \\
3  & Spectral complexity     & Spectral    & \(H_q^*\)                                                 & 5.4 \\
4  & Dissipation partitions  & Statistical & \(R_{\mathrm{grad}}^*,R_{\mathrm{pen}}^*,R_{\mathrm{part}}^*\) & 5.6 \\
5  & Energy--entropy geom.   & Statistical & \((E^*,H^*)\)                                             & 5.1, 5.6 \\
6  & Broadband cascade       & Spectral    & \(R_b^*\)                                                 & 5.4, 5.5 \\
7  & Convergence stability   & Dynamical   & \(\sigma_{\mathrm{III}}\)                                 & 5.2 \\
8  & Modal correlations      & Statistical & \(\bar{C}_{ij}\)                                          & 5.5 \\
9  & Modal overlap           & Spectral    & \(O_k^*\)                                                 & 5.4 \\
10 & Phase dynamics          & Dynamical   & \(C_{ij}^\phi\)                                           & 5.5 \\
11 & Phase--amplitude coupl. & Dynamical   & \(\rho_k\)                                                & 5.1 \\
12 & Lyapunov spectrum       & Topological & \(\lambda_i, D_{\mathrm{KY}}\)                            & 5.3 \\
13 & Attractor manifold      & Topological & \(D_{\mathrm{eff}}\)                                      & 5.1, 5.7 \\
14 & Perturbed stability     & Dynamical   & \(\sigma_\epsilon\)                                       & 5.1 \\
15 & Parameter universality  & Statistical & \(U\)                                                     & 5.7 \\
16 & Cross-consistency       & Statistical & \(C\)                                                     & 5.1--5.7 \\
\bottomrule
\end{tabular}
\end{table}


% **************************************************************************
\part{Computation}\label{part:computation}
% **************************************************************************

% ------------------------------------------------------------------
\chapter{The ED-SIM-01 Pipeline}\label{ch:pipeline}
% ------------------------------------------------------------------

\section{Architecture}

The pipeline has five layers, each reading the outputs of the layer below:

\begin{table}[ht]
\centering
\caption{Pipeline layers.}\label{tab:pipeline}
\begin{tabular}{@{}clll@{}}
\toprule
Layer & Components & Input & Output \\
\midrule
1 & Solver       & Parameters, ICs     & \texttt{time\_series.npz} \\
2 & Experiments  & Layer~1             & 64 run directories \\
3 & Invariants   & Layer~2             & Summary JSONs, figures \\
4 & Meta-analyses & Layer~3            & Synthesis JSONs \\
5 & Synthesis    & Layer~4             & Report, index, certificate \\
\bottomrule
\end{tabular}
\end{table}

\begin{figure}[ht]
\centering
% Placeholder for pipeline diagram
\fbox{\parbox{0.8\textwidth}{\centering\vspace{3cm}\textit{Five-layer pipeline diagram: Solver \(\to\) Experiments \(\to\) Invariants \(\to\) Meta-analyses \(\to\) Synthesis.}\vspace{3cm}}}
\caption{Pipeline architecture diagram. Five horizontal layers connected by vertical arrows, from the PDE solver (Layer~1) through to the certificate generation (Layer~5).}
\label{fig:pipeline}
\end{figure}


% ------------------------------------------------------------------
\chapter{Simulation Framework and Numerical Methods}\label{ch:numerics}
% ------------------------------------------------------------------

\section{Spatial Discretisation}\label{sec:spatial}

\textbf{Method~A (Finite Difference).} The spatial domain \([0,L]\) is discretised on a uniform grid of \(N+2\) points with spacing \(h = L/(N+1)\). The Laplacian is approximated by the standard three-point stencil:
\begin{equation}\label{eq:laplacian}
\Delta_h\rho_j = \frac{\rho_{j-1} - 2\rho_j + \rho_{j+1}}{h^2},
\end{equation}
with Neumann boundary conditions enforced via ghost-point reflection. The gradient squared uses central differences:
\begin{equation}\label{eq:grad-sq}
|\nabla_h\rho|^2_j = \Bigl(\frac{\rho_{j+1} - \rho_{j-1}}{2h}\Bigr)^2.
\end{equation}

\textbf{Method~B (Spectral).} The perturbation \(u = \rho - \rhostar\) is expanded in the Neumann eigenbasis~\eqref{eq:eigenbasis}. The linear part is treated exactly in spectral space; the nonlinear part is evaluated pseudospectrally with \(\tfrac{3}{2}\)-rule dealiasing.

\section{Time Stepping}\label{sec:time-stepping}

\textbf{Crank--Nicolson (default).} Diffusion is treated with half-weight implicit/explicit; nonlinear terms are explicit:
\begin{equation}\label{eq:CN}
\frac{\rho^{n+1} - \rho^n}{\Delta t} = \frac{D}{2}\bigl(L_h\rho^{n+1} + L_h\rho^n\bigr) + D\,N_h[\rho^n] - D\,P(\rho^n) + H\,v^n,
\end{equation}
where \(L_h\) is the discrete Laplacian weighted by \(M(\rho^n)\) and \(N_h[\rho^n] = M'(\rho^n)\,|\nabla_h\rho^n|^2\).

\textbf{Participation update.} The participation variable is advanced by the exact exponential integrator:
\begin{equation}\label{eq:participation-update}
v^{n+1} = e^{-\zeta\Delta t/\tau}\,v^n + \frac{\bar{F}^n}{\zeta}\bigl(1 - e^{-\zeta\Delta t/\tau}\bigr).
\end{equation}

\section{Stability Controls}\label{sec:stability}

Five structural invariants are monitored at every time step:
\begin{enumerate}
\item \textbf{Positivity}: \(\rho_j^n > 0\) for all~\(j\).
\item \textbf{Sub-capacity}: \(\rho_j^n < \rhomax\) for all~\(j\).
\item \textbf{Energy monotonicity}: \(\Energy^{n+1} \leq \Energy^n\).
\item \textbf{Mass consistency}: \(|\mathcal{M}^{n+1} - \mathcal{M}^n|\) within truncation error.
\item \textbf{Mobility positivity}: \(M(\rho_j^n) > 0\) at all interior points.
\end{enumerate}

\section{Default Numerical Parameters}\label{sec:defaults}

\begin{table}[ht]
\centering
\caption{Default numerical parameters.}\label{tab:defaults}
\begin{tabular}{@{}llp{5cm}@{}}
\toprule
Parameter & Value & Justification \\
\midrule
\(N\)           & 768                    & Converged for all Atlas experiments \\
\(\Delta t\)    & \(3.125 \times 10^{-5}\) & \(27\times\) below CFL limit \\
Method          & Crank--Nicolson        & Second-order temporal accuracy \\
\(\epsilon_{\mathrm{pos}}\) & \(10^{-12}\) & Positivity clipping threshold \\
\bottomrule
\end{tabular}
\end{table}


% ------------------------------------------------------------------
\chapter{Parameter Grid and Regime Structure}\label{ch:parameters}
% ------------------------------------------------------------------

\section{The Three-Dimensional Grid}

\begin{table}[ht]
\centering
\caption{Parameter grid for the regime volume experiment.}\label{tab:grid}
\begin{tabular}{@{}llc@{}}
\toprule
Parameter & Values & Count \\
\midrule
\(D\)   & 0.05, 0.1, 0.2, 0.5  & 4 \\
\(A\)   & 0.005, 0.01, 0.02, 0.05 & 4 \\
\(N_m\) & 5, 10, 20, 30         & 4 \\
\midrule
\multicolumn{2}{@{}l}{Total}  & 64 \\
\bottomrule
\end{tabular}
\end{table}

Each run integrates the canonical PDE from a broadband eigenmode initial condition (modes~1 through~\(N_m\) with equal amplitude~\(A\)) to final time \(T = 20\).

\section{Regime Classification}

The discriminant~\(\Disc\) (\cref{eq:discriminant}) classifies each run:
\begin{itemize}
\item \textbf{Underdamped} (\(\Disc < 0\)): the homogeneous mode oscillates before decaying.
\item \textbf{Critical} (\(|\Disc| < 10^{-10}\)): the transition between oscillatory and monotonic.
\item \textbf{Overdamped} (\(\Disc > 0\)): the homogeneous mode decays monotonically.
\item \textbf{Inadmissible}: positivity or capacity violation during integration.
\end{itemize}

\section{Five Canonical Parameter Sets}

\begin{table}[ht]
\centering
\caption{Canonical parameter sets.}\label{tab:param-sets}
\begin{tabular}{@{}cllll@{}}
\toprule
Set & \(D\) & \(\zeta\) & \(\tau\) & Regime \\
\midrule
I   & 0.3 & 0.1 & 1.0 & Deep oscillatory \\
II  & 0.6 & 0.5 & 1.0 & Moderate oscillatory \\
III & 0.8 & 1.8 & 1.0 & Near-critical \\
IV  & 0.9 & 5.0 & 1.0 & Deep monotonic \\
V   & 0.2 & 0.3 & 0.5 & High participation \\
\bottomrule
\end{tabular}
\end{table}

All sets use \(\rhostar = 0.5\), \(\rhomax = 1.0\), \(M_0 = 1.0\), \(\beta = 2.0\), \(P_0 = 1.0\).


% ------------------------------------------------------------------
\chapter{Invariant Extraction}\label{ch:extraction}
% ------------------------------------------------------------------

\section{Extraction Pipeline}

Each of the sixteen invariant scripts follows a common pattern:
\begin{enumerate}
\item \textbf{Discover} all admissible runs.
\item \textbf{Load} time-series data and metadata.
\item \textbf{Compute} the invariant over the attractor window.
\item \textbf{Fit} exponential convergence: \(|Q(t) - Q^*| \sim C\,e^{-\sigma t}\).
\item \textbf{Record} \(Q^*\), \(\sigma\), \(R^2\), and convergence flag (\(R^2 > 0.95\) and \(\sigma > 0\)).
\item \textbf{Aggregate} across runs: mean, std, min, max, CV, verdict.
\item \textbf{Produce} three standard figures.
\item \textbf{Save} summary JSON and figures.
\end{enumerate}

\section{Streaming Computation}

For memory-intensive invariants (triad balance, modal correlations), the extraction uses a streaming pattern: each triad or mode pair is processed individually, the summary statistics are recorded, and the temporary arrays are discarded. This keeps the memory footprint below 200\,MB.

\section{Convergence-Rate Fitting}

The exponential convergence rate~\(\sigma\) is extracted by linear regression of \(\log|Q(t) - Q^*|\) against~\(t\) over the fitting window. The quality metric is:
\begin{equation}\label{eq:R2}
R^2 = 1 - \frac{\sum_i (\log|Q_i - Q^*| - \hat{y}_i)^2}{\sum_i (\log|Q_i - Q^*| - \bar{y})^2}.
\end{equation}


% ------------------------------------------------------------------
\chapter{Meta-Analyses}\label{ch:meta}
% ------------------------------------------------------------------

\section{Parameter Universality}\label{sec:meta-univ}

The universality analysis constructs the pairwise distance matrix
\begin{equation}\label{eq:distance}
d_{ij} = \|\mathbf{I}_i - \mathbf{I}_j\|_2
\end{equation}
between standardised invariant vectors and computes the universality score
\begin{equation}\label{eq:U}
U = \frac{1}{1 + \CV(d)}.
\end{equation}

\section{Cross-Invariant Consistency}\label{sec:meta-consistency}

For each pair of invariant families \((A,B)\), the mean pairwise correlation is computed. The consistency score is
\begin{equation}\label{eq:C}
C = \bigl\langle |\corr(A, B)| \bigr\rangle_{A \neq B}.
\end{equation}

\begin{figure}[ht]
\centering
% Placeholder for figure
\fbox{\parbox{0.8\textwidth}{\centering\vspace{3cm}\textit{Cross-invariant correlation heatmap: \(16 \times 16\) matrix of \(|\corr(A,B)|\).}\vspace{3cm}}}
\caption{Cross-invariant correlation heatmap. Cell colour encodes \(|\corr(A,B)|\) from white~(0) through yellow~(0.5) to red~(1). Families are clustered by similarity.}
\label{fig:cross-corr}
\end{figure}

\section{Embedding Collapse}\label{sec:meta-embedding}

The standardised invariant vectors are projected into two dimensions using PCA, t-SNE, and (optionally) UMAP. The cluster radius, cluster diameter, and embedding consistency \(C_{\mathrm{emb}} = \corr(d_{\mathrm{PCA}}, d_{\mathrm{UMAP}})\) are reported.

\begin{figure}[ht]
\centering
% Placeholder for figure
\fbox{\parbox{0.8\textwidth}{\centering\vspace{3cm}\textit{Embedding map: PCA, t-SNE, and UMAP projections showing cluster collapse.}\vspace{3cm}}}
\caption{Embedding map. Three panels showing the 60~admissible runs as scatter points, coloured by~\(D\), with marker shape encoding~\(N_m\). The centroid is marked; a dashed circle indicates the cluster radius.}
\label{fig:embedding}
\end{figure}


% **************************************************************************
\part{Synthesis}\label{part:synthesis}
% **************************************************************************

% ------------------------------------------------------------------
\chapter{The Invariant Atlas}\label{ch:atlas}
% ------------------------------------------------------------------

The invariant atlas is the complete collection of sixty-four simulation runs, sixteen invariant families with per-run summaries and global statistics, three meta-analyses with synthesis figures, and a global atlas report. It covers the three-dimensional parameter space \((D, A, N_m)\) with \(4 \times 4 \times 4 = 64\) points.

The atlas is not a parameter study in the traditional sense. It does not ask ``how does observable~\(X\) depend on parameter~\(D\)?'' Instead, it asks ``is observable~\(X\) the same for all values of~\(D\)?'' The answer---quantified by the CV---determines whether \(X\) is a structural property of the architecture (\(\CV < 5\%\)) or a parameter-dependent quantity (\(\CV \geq 15\%\)).


% ------------------------------------------------------------------
\chapter{Global Verdict and Architectural Consistency}\label{ch:verdict}
% ------------------------------------------------------------------

\section{Synthesis}

\begin{table}[ht]
\centering
\caption{Global diagnostic criteria.}\label{tab:diagnostics}
\begin{tabular}{@{}llp{4cm}@{}}
\toprule
Diagnostic & Metric & Criterion \\
\midrule
Universality            & \(U\)             & \(> 0.9\): \textsc{universal} \\
Cross-Consistency       & \(C\)             & \(> 0.8\): \textsc{consistent} \\
Stability               & \(n_+, D_{\mathrm{KY}}, D_{\mathrm{eff}}\) & \(n_+ = 0\): \textsc{stable} \\
Embedding Collapse      & cluster radius    & \(< 0.1\): \textsc{collapsed} \\
Perturbation Recovery   & \(\epsilon\)-CV   & \(< 0.05\): \textsc{robust} \\
\bottomrule
\end{tabular}
\end{table}

\section{The Final Verdict}

The final verdict is:
\begin{itemize}
\item \textbf{PASS} if all five diagnostics are satisfactory.
\item \textbf{PARTIAL} if some are weakly satisfactory.
\item \textbf{FAIL} if any diagnostic contradicts the architecture.
\end{itemize}

\section{What ``Structural Consistency'' Means}

A PASS verdict does not prove that the ED architecture describes physical reality. It proves that the architecture is \emph{self-consistent}: the mathematical structure predicted by the seven principles is realised by the canonical PDE across all tested parameter regimes. The physical question---whether specific systems belong to \(\UED\)---remains empirical.


% ------------------------------------------------------------------
\chapter{The ED Architecture Certificate}\label{ch:certificate}
% ------------------------------------------------------------------

The ED Architecture Certificate is a single-page document that encodes the global verdict. It contains a header (title, version, pipeline identifier), a verdict banner (PASS/PARTIAL/FAIL with colour coding), a metrics table (five diagnostic rows with icons, labels, key metrics, and per-row verdict chips), and a footer (reproducibility statement, timestamp, version).

\begin{figure}[ht]
\centering
% Placeholder for certificate figure
\fbox{\parbox{0.8\textwidth}{\centering\vspace{4cm}\textit{ED Architecture Certificate (PASS variant). Header, verdict banner, five-row metrics table, footer, logarithmic-spiral watermark.}\vspace{4cm}}}
\caption{ED Architecture Certificate (PASS variant). The certificate summarises the global structural verdict of the ED-SIM-01 invariant atlas. Five diagnostics (universality, cross-consistency, stability, embedding collapse, perturbation recovery) are reported with per-row verdict chips. A logarithmic-spiral watermark at 3\% opacity underlies the figure.}
\label{fig:certificate}
\end{figure}


% ------------------------------------------------------------------
\chapter{Interpretation and Implications}\label{ch:implications}
% ------------------------------------------------------------------

\section{For the Architecture}

A PASS certificate confirms that the seven principles produce a coherent mathematical structure: the predictions of \cref{ch:predictions} are realised quantitatively across all tested regimes. This is the numerical foundation for the physical predictions of the Applications Paper~\cite{applications}.

\section{For the Universality Class}

The universality score~\(U\) tests the closure theorems of Appendix~D of the Rigour Paper~\cite{rigour}. A high~\(U\) confirms that the qualitative attractor structure is invariant under parameter variation within~\(\UED\).

\section{For the Physical Predictions}

The Applications Paper~\cite{applications} derives nineteen physical predictions from the architectural principles. Each prediction rests on the assumption that the target physical system belongs to~\(\UED\). The invariant atlas confirms that the canonical system realises the predicted structure; the physical question is whether specific systems can be mapped into the canonical form.


% **************************************************************************
\part{Reproducibility}\label{part:reproducibility}
% **************************************************************************

% ------------------------------------------------------------------
\chapter{Reproducibility Suite}\label{ch:repro-suite}
% ------------------------------------------------------------------

The reproducibility suite is located at \texttt{reproducibility/} within the ED~Simulation repository. The command \texttt{python reproducibility/run\_all.py} executes eight phases: environment checks, data integrity, regime volume experiment (64~runs), invariant analyses (16~families), meta-analyses (3), global atlas report, master index and certificate, and output validation.

Each phase is isolated: a failure in one phase does not abort the pipeline. The master script produces a JSON pipeline report with per-step status, timing, and error messages.

\section{Reproducibility Tiers}

\begin{table}[ht]
\centering
\caption{Reproducibility tiers.}\label{tab:tiers}
\begin{tabular}{@{}clp{5cm}@{}}
\toprule
Tier & Criterion & Scope \\
\midrule
1 & Identical file hashes & Same OS, Python, libraries, hardware \\
2 & Field-by-field agreement within \(10^{-10}\) & Same libraries, different hardware \\
3 & All invariant verdicts identical & Any conforming environment \\
\bottomrule
\end{tabular}
\end{table}


% ------------------------------------------------------------------
\chapter{Validation and Verification}\label{ch:validation}
% ------------------------------------------------------------------

The validation script \texttt{validate\_outputs.py} checks that all expected outputs exist: 64~valid regime-volume runs, 19~invariant figure directories with PNG files, 19~invariant summary JSONs, and 6~atlas files (report, index, certificate, environment, integrity, validation). It produces a JSON validation report and exits with code~0 (all checks pass) or~1 (critical failure).


% ------------------------------------------------------------------
\chapter{How to Extend the Architecture}\label{ch:extension}
% ------------------------------------------------------------------

\section{Adding a New Invariant}

\begin{enumerate}
\item Create \texttt{experiments/invariant\_my\_quantity.py} following the template.
\item Implement the standard extraction pipeline (\cref{ch:extraction}).
\item Produce the three standard figures.
\item Save a summary JSON with per-run and global statistics.
\item Add the family to the pipeline and certificate generator.
\item Run the pipeline and verify the certificate.
\end{enumerate}

\section{Porting to a New Physical Domain}

\begin{enumerate}
\item Identify the density field~\(\rho\), the constitutive functions \(M(\rho)\) and \(P(\rho)\), and the canonical parameters.
\item Verify that the seven principles (\cref{ch:principles}) are satisfied.
\item Implement the constitutive functions.
\item Run the regime volume experiment with domain-specific parameter ranges.
\item Compute the invariant atlas and compare to the canonical results.
\end{enumerate}


% **************************************************************************
\part{Outlook}\label{part:outlook}
% **************************************************************************

% ------------------------------------------------------------------
\chapter{Applications and Future Directions}\label{ch:applications}
% ------------------------------------------------------------------

The Applications Paper~\cite{applications} derives predictions for five physical domains: quantum mechanics, galactic dynamics, condensed matter, photonics, and phononics. Each prediction has a stated falsification condition. One has been tested (dwarf galaxy rotation curves) and confirmed. Eighteen additional predictions await experimental testing.

Planned extensions include higher spatial dimensions (2D, 3D), adaptive mesh refinement, GPU acceleration, and the stochastic extension:
\begin{equation}\label{eq:stochastic}
d\rho = \bigl(D\,\Fop + H\,v\bigr)\,dt + \sigma_\rho\,g(\rho)\,dW(x,t),
\end{equation}
where the mobility-proportional noise \(g(\rho) = \sqrt{M(\rho)}\) respects the capacity bound.


% ------------------------------------------------------------------
\chapter{Open Problems}\label{ch:open}
% ------------------------------------------------------------------

\section{Mathematical Open Problems}

\begin{enumerate}
\item \textbf{Multi-field extensions.} The canonical system has one density field. Multi-field generalisations and their universality class theory are open.
\item \textbf{Relativistic formulation.} A covariant formulation on curved spacetime is needed for cosmological applications.
\item \textbf{Derivation from first principles.} Whether the seven principles can be derived from a Lagrangian or action principle is unknown.
\item \textbf{Sharp universality boundaries.} The precise boundary of~\(\UED\) has not been characterised.
\end{enumerate}

\section{Computational Open Problems}

\begin{enumerate}[resume]
\item \textbf{Finer parameter grids.} The current \(4 \times 4 \times 4\) grid may miss structure near the critical surface.
\item \textbf{Long-time integration.} The standard \(T = 20\) may be insufficient for high-complexity initial conditions.
\item \textbf{Higher-dimensional invariants.} Generalisation to 2D and 3D is straightforward in principle but computationally expensive.
\end{enumerate}

\section{Physical Open Problems}

\begin{enumerate}[resume]
\item \textbf{Constitutive identification.} For each physical domain, \(M(\rho)\) and \(P(\rho)\) must be identified from the microscopic theory.
\item \textbf{Experimental falsification.} Eighteen of the nineteen predictions remain untested.
\end{enumerate}


% ------------------------------------------------------------------
\chapter{Concluding Remarks}\label{ch:conclusion}
% ------------------------------------------------------------------

This monograph has presented the Event Density architecture as a unified mathematical framework, developed its invariant theory, and documented its computational verification through the ED-SIM-01 pipeline.

The principal results are:
\begin{enumerate}
\item The canonical ED PDE possesses a unique point attractor \((\rhostar, 0)\) across all tested parameter regimes, confirmed by zero positive Lyapunov exponents and near-zero effective attractor dimension.
\item The sixteen invariant families are mutually consistent and structurally invariant under parameter variation, with coefficient of variation below 5\% for the core families.
\item The three-stage convergence structure (global bounds, algebraic decay, exponential decay) is reproduced quantitatively in every admissible run.
\item The complete pipeline---from raw PDE integration through invariant computation to certificate generation---is fully reproducible from a single command.
\end{enumerate}

The ED Architecture Certificate encodes these results in a single, machine-verifiable document. It is the bridge between the mathematical architecture and the physical predictions: if the certificate reads PASS, the architecture is self-consistent, and the predictions that follow from it are structurally well-founded.

The architecture does not explain everything. It does not fix the constitutive functions for specific physical systems. It does not produce quantitative parameter values. It does not extend to the relativistic regime. These limitations define the open problems of the programme.

What the architecture does provide is a \emph{structural guarantee}: any system satisfying the seven principles will exhibit the same qualitative dynamics---the same attractor, the same convergence, the same modal hierarchy, the same spectral fingerprint. This guarantee, once confirmed experimentally, would constitute a new kind of physical law: not a law about particles or forces, but a law about the architecture of dynamics itself.


% ==========================================================================
% APPENDICES
% ==========================================================================
\appendix

\chapter{Notation}\label{app:notation}

See the notation table in the front matter (p.~\pageref{tab:overview}).

\chapter{Key Derivations}\label{app:derivations}

\section{Energy Dissipation Identity}

Starting from~\eqref{eq:canonical-a}, multiply by \(P(\rho)/M(\rho)\) and integrate over~\(\Omega\):
\[
\int_\Omega \frac{P(\rho)}{M(\rho)}\,\partial_t\rho\,dx = D\int_\Omega \frac{P(\rho)}{M(\rho)}\,\Fop\,dx + H\,v\int_\Omega \frac{P(\rho)}{M(\rho)}\,dx.
\]
The left-hand side is \(d/dt\int_\Omega \Phi(\rho)\,dx\). After integration by parts and the participation equation, the right-hand side yields the three dissipation channels. The full derivation is in~\cite{rigour}.

\section{Triad Coupling Coefficients}

In the Neumann eigenbasis, the product \(\varphi_m'(x)\varphi_n'(x)\) projects onto \(\varphi_{m+n}\) and \(\varphi_{|m-n|}\) by the product-to-sum formula for cosines. The coupling coefficient for the forward triad \((m, n, m+n)\) is:
\begin{equation}\label{eq:Gamma}
\Gamma_{mn,m+n} = \frac{m\,n}{2\,(m+n)^2}.
\end{equation}

\section{Linearised Eigenvalue Problem}

Linearising~\eqref{eq:canonical} about \((\rhostar, 0)\) for the homogeneous mode (\(k = 0\)):
\begin{equation}\label{eq:linearised}
\begin{pmatrix} \dot{a}_0 \\ \dot{v} \end{pmatrix} = \begin{pmatrix} -D\,P_*' & H \\ -P_*'/\tau & -\zeta/\tau \end{pmatrix} \begin{pmatrix} a_0 \\ v \end{pmatrix}.
\end{equation}
The characteristic polynomial is \(\lambda^2 + 2\gamma_0\lambda + c = 0\) with \(\gamma_0 = \frac{1}{2}(D\,P_*' + \zeta/\tau)\) and \(c = (D\,P_*'\zeta + H\,P_*')/\tau\). The discriminant is \(\Disc = 4(\gamma_0^2 - c)\).


\chapter{Additional Figures}\label{app:figures}

\begin{figure}[ht]
\centering
\fbox{\parbox{0.8\textwidth}{\centering\vspace{2cm}\textit{Regime classification map: the \((D, \zeta)\)-plane partitioned into spiral, monotonic, and critical regions.}\vspace{2cm}}}
\caption{Regime classification map. The five canonical parameter sets are marked.}
\label{fig:regime-map}
\end{figure}

\begin{figure}[ht]
\centering
\fbox{\parbox{0.8\textwidth}{\centering\vspace{2cm}\textit{Three-stage convergence: semilog-\(y\) plot of \(|E(t) - E^*|\) showing Stages I, II, III.}\vspace{2cm}}}
\caption{Three-stage convergence for a representative run. Transition times between stages are marked with vertical dashed lines.}
\label{fig:three-stage}
\end{figure}


\chapter{Extended Parameter Studies}\label{app:extended}

Extended studies (not part of the standard atlas) include fine \(D\)~sweeps (\(D \in \{0.01, 0.02, \ldots, 0.99\}\)), critical-surface approaches, high-complexity frontiers (\(N_m \in \{50, 100, 200\}\)), and long-time integrations (\(T = 200\)). These extend the atlas but do not change the architectural verdict.


\chapter{Code Structure and API}\label{app:code}

\begin{table}[ht]
\centering
\caption{Core modules of ED-SIM-01.}\label{tab:modules}
\begin{tabular}{@{}lp{7cm}@{}}
\toprule
Module & Purpose \\
\midrule
\texttt{edsim\_core.py}       & PDE solver: spatial operators, time stepping, energy computation \\
\texttt{edsim\_parameters.py} & Parameter sets, constitutive functions, validation \\
\texttt{edsim\_diagnostics.py}& Observable extraction, modal decomposition, dissipation channels \\
\texttt{edsim\_initial\_conditions.py} & IC generators \\
\texttt{edsim\_runner.py}     & Integration loop, checkpointing, output logging \\
\bottomrule
\end{tabular}
\end{table}


\chapter{Glossary}\label{app:glossary}

\begin{longtable}{@{}lp{8cm}@{}}
\toprule
\textbf{Term} & \textbf{Definition} \\
\midrule
\endhead
Admissible & A density profile with \(\rho \in (0, \rhomax)\) and \(\rho \in H^1(\Omega)\) \\
Attractor & The long-time limit of all solutions; for ED, the point \((\rhostar, 0)\) \\
Attractor window & The last 10--20\% of the time series \\
Capacity bound & The maximum density~\(\rhomax\); the mobility vanishes here \\
Channel & One of the two pathways (direct or participation) \\
Constitutive function & The mobility \(M(\rho)\) or penalty \(P(\rho)\) \\
Discriminant & \(\Disc\); classifies dynamics as oscillatory, critical, or monotonic \\
Energy functional & \(\Energy[\rho,v]\); a Lyapunov functional \\
Invariant & A quantity approximately constant across parameter regimes (\(\CV < 5\%\)) \\
Invariant atlas & The complete collection of invariant computations \\
Mobility collapse & The vanishing of \(M(\rho)\) at \(\rhomax\) \\
Modal hierarchy & The ordering \(\alpha_0 < \alpha_1 < \cdots\) \\
Participation variable & \(v(t)\); integrates the spatial average of \(\Fop\) \\
Penalty & \(P(\rho)\); the restoring force toward~\(\rhostar\) \\
Regime & Classification by discriminant \\
Triad & A triple \((m,n,k)\) with \(k = m+n\) or \(k = |m-n|\) \\
Universality class & \(\UED\); the set of all systems satisfying the seven principles \\
Universality score & \(U = 1/(1 + \CV(d))\) \\
\bottomrule
\end{longtable}


% ==========================================================================
% BIBLIOGRAPHY
% ==========================================================================
\printbibliography[heading=bibintoc,title={References}]

\end{document}
